\section*{Chapter \ref{ch:g12:age-of-the-earth} --- Radiometric Dating and the Age of the Earth}

\begin{solution}{exo:g12:age-of-the-earth:1}
(a) a \omterm{def:g12:age-of-the-earth:closed-system}{closed system}; (b) an \omterm{def:g12:age-of-the-earth:isochron}{isochron}; (c) the \omterm{def:g12:age-of-the-earth:age}{age of the Earth}.
\end{solution}

\begin{solution}{exo:g12:age-of-the-earth:2}
Plagioclase (0.35), hornblende (1.20), whole rock (3.10), K-feldspar (6.40),
biotite (24.0). The biotite plots highest: it holds the most rubidium-87 for
its strontium-86, so its rubidium has made the most strontium-87 since the
granite formed.
\end{solution}

\begin{solution}{exo:g12:age-of-the-earth:3}
$T = \ln 2/\lambda = 0.6931/\num{1.3972e-11} \approx \num{4.96e10}$ years,
49.6 billion years: $49.6/4.54 \approx 11$ times the \omterm{def:g12:age-of-the-earth:age}{age of the Earth}.
\end{solution}

\begin{solution}{exo:g12:age-of-the-earth:4}
$\lambda t = \num{1.3972e-11} \times \num{3.10e8} \approx 0.004331$, and
$e^{0.004331} - 1 \approx 0.004341$. For 480 million years, $\lambda t
\approx 0.006707$ and $e^{\lambda t} - 1 \approx 0.006729$. Both are close to
$\lambda t$ itself, since $\lambda t$ is small.
\end{solution}

\begin{solution}{exo:g12:age-of-the-earth:5}
Phillips (1860), Kelvin (1862), Joly (1899), Holmes (1927), Patterson
(1956). $4540/100 \approx 45$ times.
\end{solution}

\begin{solution}{exo:g12:age-of-the-earth:6}
$0.7080 + 10.0 \times 0.004341 \approx 0.7514$ today; 0.7080 when the granite
formed, like every mineral of the granite.
\end{solution}

\begin{solution}{exo:g12:age-of-the-earth:7}
(1) $m = (0.77525 - 0.71806)/(9.40 - 0.90) = 0.05719/8.50 \approx 0.006728$;
$t = \ln(1.006728)/\num{1.3972e-11} \approx \num{4.80e8}$ years, 480 million
years. (2) $0.71806 + 1.70 \times 0.006728 \approx 0.72950$ and $0.71806 + 4.20
\times 0.006728 \approx 0.74632$: both on the line. (3) Yes: the granite cuts
the schist, so it is younger, and 310 is less than 480 million years.
\end{solution}

\begin{solution}{exo:g12:age-of-the-earth:8}
All its points would fall at the same place: one point defines no line and
no slope. The biotite lies far to the right of the others; a small error on
any point then changes the slope of the line through it and the plagioclase
much less than it would for two points close together.
\end{solution}

\begin{solution}{exo:g12:age-of-the-earth:9}
(1) Its rubidium is unchanged, so it keeps its $x$; it has less
strontium-87, so its $y$ falls: the point moves straight down. (2) The line
through it and the plagioclase is less steep, and gives an age too young.
(3) The biotite no longer lies on the line through the other minerals: the
points are not aligned, a sign that one sample has not been a \omterm{def:g12:age-of-the-earth:closed-system}{closed system}.
\end{solution}

\begin{solution}{exo:g12:age-of-the-earth:10}
(1) $4540/96 \approx 47$ times. (2) If some salt leaves the sea, for example
buried in sediments, the sea holds less salt than the rivers brought in all:
dividing the salt present by the yearly supply gives a time that is too
short.
\end{solution}

\begin{solution}{exo:g12:age-of-the-earth:11}
$m = (34.86 - 10.34)/(50.28 - 9.46) = 24.52/40.82 \approx 0.601$, between
the slopes for 4.55 (0.590) and 4.60 billion years (0.610): about 4.58
billion years.
\end{solution}

\begin{solution}{exo:g12:age-of-the-earth:12}
(1) The basalt, with the lower ratio, from the mantle; the granite, with the
higher one, at least partly from melted crust. (2) Rubidium goes into the
melt when the mantle melts, so it has been concentrated in the crust, which
has made more strontium-87 over the ages. (3) No: minerals that crystallise
from one magma share its starting ratio, and these two rocks do not.
\end{solution}

\begin{solution}{exo:g12:age-of-the-earth:13}
(1) Rubidium-87: $2^{-4.54/49.61} \approx 0.94$, about \qty{94}{\percent}
left; uranium-238: $2^{-4.54/4.47} \approx 0.49$, about half. (2) $m =
e^{\lambda t} - 1 \approx \num{1.3972e-11} \times 10^6 \approx
\num{1.4e-5}$: even the biotite of the Varden granite ($x = 24$) would rise
by only $0.0003$ above the starting ratio. So little strontium-87 has formed
that the slope is hard to tell from zero: a slow clock suits old rocks.
\end{solution}

\begin{solution}{exo:g12:age-of-the-earth:14}
(1) $4.438 + 0.6026 \times 19.0 \approx 15.89$; $15.5$ is lower by $0.39$,
about \qty{2.4}{\percent}: the sediment lies close to the line. (2) 0.6026
lies between 0.590 (4.55) and 0.610 (4.60): about 4.58 billion years, within
Patterson's range of 4.48 to 4.62 billion years.
\end{solution}

\begin{solution}{exo:g12:age-of-the-earth:15}
(1) The store then loses only $46.7 - 24.4 = \qty{22.3}{TW}$ net, so the same
store lasts $46.7/22.3 \approx 2.1$ times longer. (2) $96 \times 2.1 \approx
200$ million years, still more than twenty times too short. (3) \omterm{def:g12:internal-heat:convection}{Mantle
convection}, which carries heat up from deep inside and keeps the gradient
near the surface high far longer than conduction alone.
\end{solution}

\begin{solution}{pb:g12:age-of-the-earth:1}
\textbf{1.} The biotite ($x = 24.0$).

\textbf{2.} They crystallised from one magma, and minerals that crystallise
from one melt take in its strontium with its ratio of isotopes.

\textbf{3.} Strontium-86 is stable and made by no decay: in a \omterm{def:g12:age-of-the-earth:closed-system}{closed system}
its number does not change, so it is a fixed yardstick.

\textbf{4.} Strontium-87, made by the decay of rubidium-87; rubidium-87 has
decreased and strontium-86 has stayed the same.

\textbf{5.} A horizontal line at $y = 0.7080$: same ratio, different $x$.

\textbf{6.} $(0.81218 - 0.70952)/(24.0 - 0.35) = 0.10266/23.65 \approx
0.004341$.

\textbf{7.} $(0.72146 - 0.70952)/(3.10 - 0.35) = 0.01194/2.75 \approx
0.004342$.

\textbf{8.} They agree to within 1 in 4000. The first is the more reliable:
its two points are far apart, so the rounding of the ratios matters less.

\textbf{9.} $0.70952 - 0.35 \times 0.004341 \approx 0.7080$.

\textbf{10.} $0.7080 + 6.40 \times 0.004341 \approx 0.73578$: exactly the
measured value.

\textbf{11.} $e^{\lambda t} = 1 + m = 1.004341$.

\textbf{12.} $\lambda t = \ln 1.004341 \approx 0.004331$.

\textbf{13.} $t = 0.004331/\num{1.3972e-11} \approx \num{3.10e8}$ years.

\textbf{14.} Yes: the granite cuts the schist, so by the principle of
cross-cutting relationships it is younger, and 310 is less than 480 million
years.

\textbf{15.} 310 million years lies between 358.9 and 298.9 million years:
the Carboniferous.

\textbf{16.} The ratio $^{87}\mathrm{Sr}/^{86}\mathrm{Sr}$ of the magma when
the granite crystallised.

\textbf{17.} The granite: $0.7080 > 0.7035$.

\textbf{18.} Rubidium goes into the melt when the mantle melts, so it has
been concentrated in the crust; with more rubidium, the crust has made more
strontium-87.

\textbf{19.} That the magma came, at least in part, from melted \omterm{def:g12:continental-crust:continental-crust}{continental
crust} rather than from the mantle alone.

\textbf{20.} The Varden granite is \textbf{310 million years} old.
\end{solution}
